Second derivative of free energy with respect to chromium composition. Spinodal demixing is stable between 0.26$\pm$0.01 and 0.95$\pm$0.02 chromium content. \cite{korzhavyi2009electronic,trochet2021influence,levesque2011simple, binder2001spinodal}


\subsection{Minimizing elastic anisotropy through spinodal demixing}
Based on Kádas et al. (Solid State Communications, 2011) and the European Research Council's Final Report Summary (Atomic-Level Physics of Advanced Materials, 2014), we explore the impact of spinodal demixing on elastic properties and material stability. Based on these calculations, spinodal demixing appears to play a significant role in reducing elastic anisotropy induced at spinodal bounds to minimize free energy. 

Polycrystalline anisotropy of body-centered cubic iron-chrome as a function of increasing chromium content. Increased alloying decreases elastic anisotropy between spinodal bounds to dynamically stabilize alloy. \cite{kadas2011thermo,euroresearchcouncil2014}

Pugh's classical criterion and Pettifor's quantum criterion are discussed as complementary approaches to characterizing the ductile-brittle transition. The proposed method for investigating demixing and embrittlement phenomena is scalable to high-throughput datasets. Based on Senkov and Miracle (Scientific Reports, 2021), we highlight the potential for applying this method to large-scale materials datasets. The inclusion of validation measurements ensures the reliability and accuracy of the results.


